\documentclass[10pt,a4paper]{amsart}
\usepackage[margin=19mm]{geometry}
\usepackage{amsmath,amssymb,mathtools}
\usepackage[scr=rsfs,cal=euler]{mathalfa}
\usepackage{xfrac}
\usepackage[unicode,hidelinks,bookmarks=false]{hyperref}
\newcommand{\ie}{i.e.\ }
\newcommand{\cf}{cf.\ }
\newcommand{\resp}{respectively\ }
\newcommand{\Subsection}{\subsection}
\providecommand{\nobreakdash}{\nobreak\hbox{-}}
\setlength{\emergencystretch}{2em}
\multlinegap0pt
\usepackage{bm}
\usepackage{bbm}
\usepackage{dsfont}
\usepackage{xspace}
\usepackage{cancel}
\usepackage{mathtools}
\usepackage{amsthm}
\makeatletter
\@ifundefined{thm}{\newtheorem{thm}{Thm}}{}
\@ifundefined{lem}{\newtheorem{lem}[thm]{Lemma}}{}
\makeatother
\makeatletter
\def\Z{\mathbb{Z}}
\def\a{\alpha}
\def\b{\beta}
\def\cS{\mathcal{S}}
\def\cT{\mathcal{T}}
\expandafter\ifx\csname pf\endcsname\relax
\newenvironment{pf}{\par\medskip\noindent{\bf Proof.}}{\hfill$\square$\par\medskip}
\fi
\expandafter\ifx\csname spf\endcsname\relax
\newenvironment{spf}{\par\medskip\noindent{\bf Sketch Proof.}}{\hfill$\square$\par\medskip}
\fi
\def\wh{\widehat}
\makeatother
\title[Generalising Collins' Theorem]{Generalising Collins' Theorem}
\author{James Howie, Hamish Short}
\date{Source: arXiv:2307.15397v2 (CC BY 4.0)}
\begin{document}
\maketitle

\noindent\emph{Source and licence.} Generalising Collins' Theorem. Version arXiv:2307.15397v2 (CC BY 4.0), distributed under the Creative Commons Attribution 4.0 International licence, as recorded in the arXiv licence metadata for this exact version and shown on the abstract page https://arxiv.org/abs/2307.15397v2. Full rights notice: licenses/arxiv-2307.15397-CC-BY-4.0.txt.\par\smallskip

\noindent\textit{Editorial scope.} Counted unit: the Lem whose source label is \texttt{surlepont} in the article (source0.tex lines 1301--1306, source label \texttt{surlepont}), delivered together with the article's own complete original proof of it (source0.tex lines 1308--1347, one \texttt{proof} environment of 40 lines). No mathematics is written, repaired or re-derived: statement and proof are the article's own text line for line. Required context retained uncounted: none. Every definition, display and label of those lines is retained unchanged. Omitted from this package and not redistributed: 1--1300, 1307--1307, 1348--1686. Nothing inside a retained excerpt is omitted or altered: every retained body is delivered exactly as the article's lines are written, so no omission receipt of any kind accompanies this package. Citations: the retained excerpts carry no citation command at all, so no bibliography entry of the article is transcribed and no reference is redistributed. Reference adaptation: none was needed and none was made; every reference inside the delivered unit resolves inside it, so no unresolved reference or citation is produced. Printed slips kept literal: no printed word, symbol, hypothesis or number of the counted statement or of its proof was altered. Layout and class: the article's own class is \texttt{amsart} with packages including amssymb, graphicx, tikz-cd, mathtools; the wrapper is the lane's pinned \texttt{amsart} editorial wrapper at 10pt on A4, which restates the theorem-like environments the retained text needs and adds the article's own macro definitions that the retained text needs (\texttt{\textbackslash Z}, \texttt{\textbackslash a}, \texttt{\textbackslash b}, \texttt{\textbackslash cS}, \texttt{\textbackslash cT}, \texttt{\textbackslash wh}), with extra wrapper packages (bm, bbm, dsfont, xspace, cancel, mathtools). The delivered numbering is the unit's own. Assets: no retained span contains a figure environment, a graphics-inclusion command, a PDF-page inclusion, a table or a TikZ picture; no image file is copied into this package, the package declares no asset dependency and no image file is redistributed with it. Licence: version arXiv:2307.15397v2 of this article is distributed under the Creative Commons Attribution 4.0 International licence, as recorded in the arXiv licence metadata for this exact version and shown on the abstract page https://arxiv.org/abs/2307.15397v2. A full rights notice is delivered at licenses/arxiv-2307.15397-CC-BY-4.0.txt. Characters: every retained excerpt is pure ASCII, so the delivered engine, XeLaTeX with its default Latin Modern text font, reports no missing character.
 \begin{lem}	\label{surlepont}
 	Let $P$ be a reduced annular relative picture over $\wh Y$, 
 	such that the $\a$-discs of $P$ that are joined by $\cS$-arcs 
 	to the top boundary have two or more distinct $\Z$-labels.  
 	Then $P$ has at least two up-down connections, one $G$-max and one   $G$-min.
 \end{lem}

 \begin{proof}
 	Suppose that $\b_0,\b_1$ are two $\a$-discs of $P$ that are attached to the top boundary by $\cS$-arcs, and have $\Z$-labels $n_0,n_1$ respectively with $n_0<n_1$.   
 	If $\b_1$ is not a $G$-max up-down connection, 
 	then for some $u\in\cS\cup\cT$  its $u_{max}$-arc connects $\b_1$ to another $\a$-disc -- say $\b_2$.    
 	Iterating this process gives a  chain 
  $\b_1,\b_2,\cdots$ of $\a$-discs in $P$,  
 	where each $\b_j$ is joined to $\b_{j+1}$ by the  $u_{max}$-arc  of $\b_j$ for some $u\in\cS\cup\cT$.   
 	The chain ends in a  $G$-max up-down connection, 
 	or it contains a repetition -- say $\b_i=\b_j$ with $1\le i<j$. 
 	In the latter case, 
 		assuming that $\b_i=\b_j$ is the first repetition in the chain,
 	the path $\b_i$ -- $\b_{i+1}$ -- $\cdots$ -- $\b_j=\b_i$ 
 	 is simple, and cannot be nullhomotopic in 
 	 the ambient annulus, 
 	else there  would be a 
 	 disc-subpicture containing the path and the pre-ordering in this  disc-picture
 	would give a strictly $G$-increasing chain of $\a$ discs $\b_i<\b_{i+1}<\cdots<\b_j$ which 
 		could not be closed.   
 	 Hence
 	this path
 	 $\b_i$ -- $\b_{i+1}$ -- $\cdots$ -- $\b_j=\b_i$
 		wraps once around the annulus and cuts it 
 	into two smaller annuli.  
 	Then the path $\b_i$ -- $\b_{i-1}$ -- $\cdots$ -- $\b_1$ -- [top boundary] 
	splits the upper small annulus into a disc --  $D$ say.
 	
 	Now repeat this argument in the other direction from $\b_0$ to construct a 
 	chain $\b_0$, $\b_{-1}$, \dots with 
 	$\b_{-j}$ joined to $\b_{-j-1}$ 
 	by the $u_{min}$-arc of $\b_{-j}$ for some $u\in\cS\cup\cT$.  
 	The two chains cannot meet as the $\Z$ labels in the first are all strictly 
 	greater than the $\Z$ labels in the second
  -- by our choice of right-ordering $<$.
 	The second chain is thus entirely contained in the disc $D$, 
 	so cannot contain a repetition or an up-down connection, so continues indefinitely, 
 	contradicting the fact that $P$ has only finitely many $\a$-discs.  
 	Hence the first chain $\b_1,\b_2,\cdots$ above must end in a $G$-max up-down connection.   
 	For similar reasons, the second chain  $\b_0,\b_{-1}\cdots$ must end in a  
 	$G$-min up-down connection.
 \end{proof}

\end{document}
